\problema{Minimum Window Substring}{76}{src/06-sliding-window-stack/76-minimum-window-substring.cpp}{H}

Dadas dos cadenas \texttt{s} y \texttt{t}, devuelve la subcadena
\emph{contigua} más corta de \texttt{s} que contiene a todos los caracteres
de \texttt{t}, contando repeticiones (si \texttt{t} tiene dos \texttt{'a'},
la ventana debe tener al menos dos \texttt{'a'}). Si no existe tal
subcadena, devuelve la cadena vacía.

\paragraph{La idea, con un ejemplo de la vida real.} Imagina que tienes una
lista de compras (\texttt{t}) y caminas por los pasillos de un supermercado
muy largo (\texttt{s}) tomando productos según los vas encontrando, sin
poder saltarte pasillos ni regresar por productos sueltos: solo puedes
quedarte con un tramo continuo de pasillo. Quieres encontrar el tramo más
corto que, entre todo lo que recogiste ahí, cubra tu lista completa (aunque
recojas cosas de más que no estaban en la lista, no importa). En cuanto tu
carrito ya cubre la lista completa, empiezas a soltar productos por el lado
donde entraste primero, mientras la lista siga cubierta, para ver si el
tramo se puede acortar todavía más.

\begin{center}
\ttfamily
\begin{tabular}{l}
A D O B E C O D E B A N C\\
$\phantom{AAAAAAAAAAAAAAAAAA}$[B A N C]\\
\end{tabular}
\end{center}

\noindent Con \texttt{s = "ADOBECODEBANC"}, \texttt{t = "ABC"}, la ventana
más corta que cubre \texttt{A}, \texttt{B} y \texttt{C} es \texttt{"BANC"},
de longitud \texttt{4}.

\paragraph{Cómo lo resuelve el código, paso a paso.} Guardamos en un mapa
\texttt{necesita} cuántas veces aparece cada carácter en \texttt{t}, y en
\texttt{requeridos} cuántos caracteres \emph{distintos} hay que cubrir.
Llevamos otro mapa \texttt{ventana} con lo que hay actualmente entre los
punteros \texttt{L} y \texttt{R}, y un contador \texttt{cubiertos} de
cuántos caracteres distintos ya alcanzaron la cantidad necesaria.
\begin{enumerate}
  \item Avanzamos \texttt{R} y metemos \texttt{s[R]} a la ventana.
  \item Si \texttt{s[R]} es uno de los caracteres necesarios y la ventana
        acaba de alcanzar justo la cantidad requerida de él, sumamos uno a
        \texttt{cubiertos}.
  \item Mientras \texttt{cubiertos} sea igual a \texttt{requeridos} (la
        ventana ya cubre todo \texttt{t}): comparamos su tamaño contra la
        mejor guardada, y si es más chica la guardamos; luego sacamos
        \texttt{s[L]} de la ventana y avanzamos \texttt{L}. Si al sacar
        \texttt{s[L]} la cantidad cae por debajo de lo necesario,
        restamos uno a \texttt{cubiertos} (la ventana deja de estar
        completa y el ciclo de contracción se detiene).
  \item Repetimos hasta que \texttt{R} llega al final de \texttt{s}.
\end{enumerate}
La clave es que \texttt{L} y \texttt{R} solo avanzan hacia adelante, nunca
retroceden: por eso cada carácter de \texttt{s} entra y sale de la ventana
a lo más una vez, y el algoritmo entero es lineal.

\lstinputlisting{src/06-sliding-window-stack/76-minimum-window-substring.cpp}

\complejidad{$O(|s| + |t|)$}{$O(|s| + |t|)$}

\paragraph{¿Qué significa esa complejidad?} Recorremos \texttt{s} una vez
para expandir y, en total, otra vez más para contraer (nunca más de
\texttt{|s|} pasos combinados), y construir el mapa de \texttt{t} cuesta
\texttt{|t|}. El espacio lo usan los dos mapas de conteo, acotados por
cuántos caracteres distintos puede haber.

\paragraph{Consejos para la entrevista (Amazon).} Es uno de los problemas de
sliding window más pedidos porque combina ventana de tamaño variable con
conteo de multiplicidad (no basta con ``¿está el carácter?'', hay que
contar cuántas veces se necesita). El error más común es comparar
\texttt{ventana[c] > 0} en vez de \texttt{ventana[c] >= necesita[c]}, lo
cual falla en cuanto \texttt{t} tiene caracteres repetidos. Otro punto que
suelen evaluar: por qué la condición del \texttt{while} de contracción usa
igualdad exacta (\texttt{cubiertos == requeridos}) y no una comparación de
mayor o igual —conviene explicar que \texttt{cubiertos} nunca puede superar
\texttt{requeridos} por construcción—. Casos límite a mencionar: \texttt{s}
o \texttt{t} vacías, y \texttt{t} pidiendo más repeticiones de un carácter
de las que existen en \texttt{s} (no hay solución, se devuelve cadena
vacía).
